% Part: normal-modal-logic
% Chapter: frame-correspondence
% Section: first-order-definability

\documentclass[../../../include/open-logic-section]{subfiles}

\begin{document}

\olfileid{nml}{frd}{fol}

\olsection{د لومړۍ درجې تعريف‌وړتيا}

مونږ وليدل چې د چوکاټونو د لاسرسي اړيکو يو شمېر خاصيتونه د مودالي
!!{formula}s په واسطه تعريفېدے شي۔ د بېلګې په توګه، د چوکاټونو
تناظر د~\Ax{B} په !!{formula}، $p \lif \Box\Diamond
p$، تعريفېدے شي۔ تر اوسه چې کوم شرطونه مو ليدلي، ټول د لومړۍ درجې
په هغو !!{formula}s بيانېدے شي چې په داسې !!a{language} کښې وي
چې يوازې يو دوه‌ځاييز !!{predicate} لري۔ د بېلګې په توګه، تناظر
د~$\lforall[x][\lforall[y][(\Atom{Q}{x,y} \lif \Atom{Q}{y, x})]]$
په واسطه تعريفېږي: د لومړۍ درجې هغه !!{structure}~$\Struct{M}$
چې $\Domain{M} = W$ او $\Assign{Q}{M} = R$ وي، پورته فارمول
هله او يوازې هله پوره کوي چې $R$ متناظره وي۔ دا لاندې تعريف ته
لار هواروي:

\begin{defn}
  د چوکاټونو يو ټولګے~$\mClass{F}$ \emph{د لومړۍ درجې له مخې
  تعريف‌وړ} دے، که د لومړۍ درجې په داسې ژبه کښې چې يوازې يو
  دوه‌ځاييز !!{predicate}~$Q$ لري، يوه !!a{sentence}~$!A$ وي،
  داسې چې $\mModel{F} =
  \tuple{W, R} \in \mClass{F}$ هله او يوازې هله چې د لومړۍ درجې
  په هغه !!{structure}~$\Struct{M}$ کښې $\Sat{M}{!A}$ وي چې
  $\Domain{M} = W$ او $\Assign{Q}{M}
  = R$ لري۔
\end{defn}

تر اوسه څېړل شوي خاصيتونه او هغه مودالي !!{formula}s چې تعريفوي
ئې، استثنايي دي۔ هر !!{formula} د چوکاټونو داسې ټولګے نۀ
تعريفوي چې د لومړۍ درجې له مخې تعريف‌وړ وي؛ او د چوکاټونو هر
د لومړۍ درجې له مخې تعريف‌وړ ټولګے هم په مودالي فارمول نۀ
تعريفېږي۔

د لومړۍ دعوې مقابله بېلګه د لوب فارمول دے:
\begin{equation}
\Box(\Box p \lif p) \lif \Box p. \tag{\Ax{W}}
\end{equation}
\Ax{W} د هغو چوکاټونو ټولګے تعريفوي چې اړيکه ئې انتقالي او
سرچپه ښه‌بنسټه وي۔ يوه اړيکه هله ښه‌بنسټه ده چې داسې هېڅ
نامتناهي لړۍ $w_1$، $w_2$، \dots{} ونۀ لري چې $Rw_2w_1$،
$Rw_3w_2$، \dots وي۔ د بېلګې په توګه، پر~$\Nat$ د $<$ اړيکه
ښه‌بنسټه ده، خو پر~$\Int$ د $<$ اړيکه داسې نۀ ده۔ يوه اړيکه هله
او يوازې هله سرچپه ښه‌بنسټه ده چې د هغې سرچپه اړيکه
ښه‌بنسټه وي۔ نو سرچپه ښه‌بنسټې اړيکې هغه دي چې هېڅ داسې نامتناهي
لړۍ $w_1$، $w_2$، \dots{} ونۀ لري چې $Rw_1w_2$،
$Rw_2w_3$، \dots وي۔

خو د لومړۍ درجې هېڅ !!{formula} انتقالي او سرچپه ښه‌بنسټې
اړيکې نۀ تعريفوي۔ د دې د ليدلو دپاره فرض کړئ چې
$\Sat{M}{!F}$ هله او يوازې هله وي چې $R =
\Assign{Q}{M}$ انتقالي او سرچپه ښه‌بنسټه وي۔ $!A_n$ دا
!!{formula} ونيسئ
\[
(\Atom{Q}{a_1,a_2} \land \dots \land
    \Atom{Q}{a_{n-1},a_{n}})
\]
دلته د شمېر لرونکي ثابت سمبولونه د لومړۍ درجې د ژبې په موقتي
پراختيا کښې کاروو؛ د لومړي شاخص په مورد کښې عطف تش او ځکه رښتيا
ګڼل کېږي۔ \emph{د سرچينې يادونه: اصلي ننداره د لومړي شاخص
تش عطف او د نويو ثابتونو د ژبې پراختيا په څرګنده نۀ بيانوي۔}
اوس د !!{formula}s دا سټ وڅېړئ
\[
\Gamma = \{!F, !A_1, !A_2, \dots\}.
\]
د~$\Gamma$ هر متناهي فرعي سټ د صدق وړ دے: $k$ هغه تر ټولو لوے
شاخص ونيسئ چې $!A_k$ په فرعي سټ کښې وي؛ که د لړۍ هېڅ غړے پکې
نۀ وي، لومړے شاخص ونيسئ۔ بيا $\Domain{M_k} = \{1, \dots, k\}$،
$\Assign{a_i}{M_k} = i$ د همدغه حد پورې، او
$\Assign{Q}{M_k} = <$ وټاکئ۔ له دې حد څخه د پورته ثابتونو ارزښتونه
په دامنه کښې خپلسري وټاکئ۔ ځکه چې پر $\{1,
\dots, k\}$ د $<$ اړيکه انتقالي او سرچپه ښه‌بنسټه ده،
$\Sat{M_k}{!F}$۔ د جوړښت له مخې، د هر $i \le k$ دپاره
$\Sat{M_k}{!A_i}$۔ \emph{د سرچينې يادونه: اصلي ثبوت د داسې
فرعي سټ حالت نۀ يادوي چې د لړۍ هېڅ غړے پکې نۀ وي، او د
ثابتونو ټاکنه يوازې تر ياد شوي حد پورې په همدې بڼه شونې ده۔}
د لومړۍ درجې منطق د متناهي شاهد قضيې له مخې، $\Gamma$ په
کوم !!{structure}~$\Struct{M}$ کښې د صدق وړ دے۔ د فرض له مخې،
ځکه چې $\Sat{M}{!F}$، نو $\Assign{Q}{M}$ سرچپه
ښه‌بنسټه ده۔ خو څرګنده ده چې $\Assign{a_1}{M}$،
$\Assign{a_2}{M}$، \dots{} به هغه نامتناهي لړۍ جوړه کړي
چې سرچپه ښه‌بنسټوالی ئې ردوي۔

د دويمې دعوې مقابله بېلګه د کليت خاصيت دے: د هر $u$ او $v$
دپاره $Ruv$۔ کلي چوکاټونه د لومړۍ درجې د
$\lforall[x][\lforall[y][\Atom{Q}{x,y}]]$ فارمول په واسطه
تعريفېدے شي۔ خو داسې هېڅ مودالي فارمول نشته چې په ټولو او يوازې
کلي چوکاټونو کښې معتبر وي۔ دا د يوې په خپل ذات کښې په زړه پورې
پايلې نتيجه ده: په کلي چوکاټونو کښې معتبر فارمولونه عين هغه
فارمولونه دي چې په انعکاسي، متناظرو او انتقالي چوکاټونو کښې
معتبر دي۔ ځينې انعکاسي، متناظره او انتقالي چوکاټونه کلي
نۀ دي؛ نو هر !!{formula} چې په ټولو کلي چوکاټونو کښې معتبر وي،
په ځينو ناکلي چوکاټونو کښې هم معتبر وي۔

\end{document}
